\documentclass[11pt,a4paper]{article}

\usepackage[T1]{fontenc}
\usepackage{lmodern}
\usepackage[utf8]{inputenc}
\usepackage[spanish,es-noquoting]{babel}
\usepackage[margin=2.6cm]{geometry}
\usepackage{amsmath,amssymb}
\usepackage{booktabs}
\usepackage{graphicx}
\usepackage{microtype}
\usepackage{enumitem}
\usepackage[table]{xcolor}
\usepackage{titlesec}
\usepackage[hidelinks]{hyperref}

% babel-spanish redefine \% para insertar un espacio fino antes del signo, y esa
% redefinicion usa aritmetica de \lastskip que rompe dentro de minipage.
\AtBeginDocument{\def\%{\char37\relax}}

\definecolor{acento}{HTML}{1F4E79}
\definecolor{suave}{HTML}{F2F5F8}

\titleformat{\section}{\large\bfseries\color{acento}}{\thesection}{0.7em}{}
\titleformat{\subsection}{\normalsize\bfseries}{\thesubsection}{0.6em}{}

\newcommand{\term}[1]{\textbf{#1}}
\setlist[itemize]{leftmargin=1.2em,itemsep=2pt,topsep=3pt}

\title{\vspace{-1.5cm}\textbf{Gray-box Wilson-Cowan: del plan del 10-09 a las
primeras correcciones que funcionan}}
\author{Informe de trabajo}
\date{11 de septiembre de 2026}

\begin{document}
\maketitle
\vspace{-0.6cm}

\noindent\fcolorbox{acento!30}{suave}{%
\begin{minipage}{\dimexpr\textwidth-2\fboxsep-2\fboxrule}
\small
El plan del 10 de septiembre afirmaba que la corrección neuronal fracasaba por
falta de información y no de capacidad, y que darle la historia del comando la
haría funcionar. Se implementó y se midió. El NRMSE sobre la planta con
actuador bajó de $15{,}23\,\%$ a $10{,}16\,\%$ y el $R^2$ de la corrección
rompió el techo teórico que ninguna función del estado observado podía superar.
Quedan corriendo doce experimentos más. Este documento recorre lo hecho desde
la escritura del plan, explica cada término y el motivo de cada decisión.
\end{minipage}}

\section{De qué trata el proyecto}

El objetivo es identificar la dinámica de una población neuronal a partir de
datos, combinando una ecuación conocida con una parte aprendida. La ecuación
conocida es el \term{modelo de Wilson-Cowan}, un modelo de masa neuronal que
describe dos poblaciones acopladas, una excitatoria y una inhibitoria, mediante
dos ecuaciones diferenciales ordinarias. Sus variables de estado son $E$, la
actividad de la población excitatoria, e $I$, la de la inhibitoria. El estado
del sistema es el vector $x = (I, E)^{\mathsf{T}}$. Las entradas externas son
$P$ y $Q$, los estímulos que recibe cada población.

El modelo tiene diez parámetros que se agrupan en el vector
$\beta \in \mathbb{R}^{10}$: cuatro pesos de acoplamiento
($w_{EE}$, $w_{EI}$, $w_{IE}$, $w_{II}$), dos constantes de tiempo
($\tau_e$, $\tau_i$), dos pendientes de la función de activación
($a_e$, $a_i$) y dos umbrales ($\theta_e$, $\theta_i$). El nombre $\beta$ se
eligió porque la letra $\theta$, la habitual para un vector de parámetros,
colisiona con los umbrales $\theta_e$ y $\theta_i$ del propio modelo.

\subsection{Qué es un gray-box y por qué no un black-box}

Un \term{black-box} aprende la dinámica entera con una red neuronal, sin
suponer nada. Reproduce bien y no explica nada: los pesos aprendidos no
corresponden a ninguna cantidad física. Un \term{white-box} escribe la ecuación
completa a mano y ajusta solo sus parámetros. Explica todo lo que acierta, pero
si la planta real tiene algo que la ecuación no contempla, ese algo no se puede
representar y el error queda.

El \term{gray-box} es la combinación:
\begin{equation}
  \dot{x} = \underbrace{f_{WC}(x, P, Q; \beta)}_{\text{física escrita a mano}}
          + \underbrace{g_\varphi(\cdot)}_{\text{red neuronal}}
\end{equation}
El primer término aporta la estructura conocida y sus parámetros conservan
significado físico. El segundo, con parámetros $\varphi$, absorbe lo que la
estructura no cubre. La promesa es tener las dos cosas: reproducción de un
black-box e interpretabilidad de un white-box.

El problema central del proyecto es que esa promesa venía sin cumplirse. En las
dos únicas configuraciones que reproducían bien la dinámica, la red estaba
literalmente apagada (\texttt{use\_correction=False}) y todo el trabajo lo hacía
física escrita a mano. Eso es lo que disparó el plan del 10 de septiembre.

\subsection{La planta simulada}

Los datos son de simulador, no experimentales. Se genera una trayectoria con
Wilson-Cowan más una \term{perturbación} que el modelo reducido no contempla, y
después se pide identificar el sistema usando solo el modelo reducido más la
corrección. La diferencia entre la dinámica verdadera y la que el modelo
reducido puede representar se llama $\Delta f$, y es exactamente lo que $g$
tendría que aprender.

Se usan dos plantas, cada una con una sola perturbación activa, para atacar la
dificultad de a una por vez:

\begin{itemize}
\item \texttt{refrac1}: \term{refractariedad}. Después de disparar, una neurona
  no puede volver a hacerlo de inmediato. Se modela con un factor
  $(1 - r\,x)$ que reduce la respuesta cuando la actividad es alta. Es una
  función del estado actual y de nada más.
\item \texttt{act1}: \term{actuador}. El estímulo que llega a la población no es
  el comando $u$ que se ordena, sino una versión filtrada y saturada de él. En
  concreto un \term{sistema de Wiener}, que es un filtro lineal seguido de una
  no linealidad sin memoria:
  \begin{equation}
    \frac{dP_{\text{lag}}}{dt} = \frac{P - P_{\text{lag}}}{\tau},
    \qquad \tau = 1~\text{ms},
    \qquad P_{\text{efectivo}} = A \tanh\!\left(\frac{P_{\text{lag}}}{A}\right),
    \qquad A = 3.
  \end{equation}
  Acá está la dificultad: el estímulo efectivo depende del \emph{pasado} del
  comando, no solo de su valor actual.
\end{itemize}

\subsection{Las métricas y qué mide cada una}

\begin{itemize}
\item \term{NRMSE de corrida libre} (\emph{free-run rollout}). Se le da al
  modelo identificado el estado inicial y la secuencia de estímulos, y se lo
  deja integrar la trayectoria entera sin volver a mirar los datos. Se compara
  contra la trayectoria verdadera y se normaliza. Mide reproducción de la
  dinámica. Es la métrica principal.
\item \term{$R^2$ de la corrección}. Compara, muestra a muestra, lo que la
  corrección aprendida $\hat{g}$ produce contra el término omitido verdadero
  $\Delta f$:
  \begin{equation}
    R^2 = 1 - \frac{\sum \lVert \hat{g} - \Delta f \rVert^2}
                   {\sum \lVert \Delta f - \overline{\Delta f} \rVert^2}
  \end{equation}
  Mide recuperación de física, no reproducción. Un modelo puede copiar la
  trayectoria con $R^2$ pésimo, si la corrección compensa por un camino
  distinto del verdadero.
\item \term{Error de parámetros}. Distancia relativa entre el $\beta$
  identificado y el $\beta$ verdadero del simulador. Es el control que dice si
  la corrección está deformando la física en vez de complementarla.
\item \term{Fracción de redundancia}. Cuánto de lo que hace la corrección lo
  podría haber hecho un cambio de parámetros. El punto de comparación no es
  cero: el $\Delta f$ verdadero ya puntúa $0{,}67$, y una red con pesos
  aleatorios puntúa $0{,}47$ solo por ser una función suave del estado.
\end{itemize}

\subsection{Cómo se entrena}

Se usa \term{multiple shooting}: la trayectoria se parte en ventanas cortas y
cada ventana se integra desde el estado observado en su comienzo, en vez de
integrar los 4000 pasos de un tirón. El motivo es de condicionamiento numérico.
Integrar una trayectoria larga y derivar respecto de los parámetros hace que el
gradiente explote o se desvanezca, porque el error se propaga multiplicativamente
a lo largo de la integración. Reiniciando cada 100 pasos el problema queda mucho
mejor planteado.

Después de la fase de descenso por gradiente viene una fase de \term{L-BFGS}, un
método de segundo orden que converge mucho más rápido cerca del mínimo pero que
necesita arrancar razonablemente cerca.

\section{El plan del 10-09 y su razonamiento}

\subsection{El hallazgo que lo disparó}

Preparando el póster de la Jornada se verificó en el código que en las dos
correcciones que funcionaban, la red estaba apagada. La primera,
\emph{forma exacta}, agrega la refractariedad al backbone como tres parámetros
físicos más. La segunda, \emph{estado de filtro}, aumenta el estado con dos
variables que siguen la dinámica del actuador, y son doce parámetros en total,
sin ninguna red. Las dos reproducen bien. Ninguna de las dos usa $g$.

Eso vacía el argumento del gray-box. Si lo único que funciona es escribir la
ecuación a mano, el proyecto es un white-box con pasos extra.

\subsection{El teorema y qué dice exactamente}

El \term{Teorema de Aproximación Universal} dice que una red neuronal
suficientemente ancha puede aproximar cualquier función continua \emph{de sus
entradas}, con la precisión que se quiera. La palabra que hace todo el trabajo
es \emph{de sus entradas}.

Sobre \texttt{act1}, $\Delta f$ depende del pasado del comando. Una red que solo
ve $(I, E)$ no puede representarla, porque $\Delta f$ \emph{no es una función}
de $(I,E)$: para un mismo estado observado, $\Delta f$ toma valores distintos
según lo que haya pasado antes. Esto no es una limitación de capacidad. Es una
limitación de información, y ninguna cantidad de capas la arregla.

Está medido. El \term{oráculo}, que es la mejor función posible de $(I,E)$
ajustada directamente contra $\Delta f$ sin ninguna restricción de
optimización, llega a $R^2 = -0{,}11$ sobre la planta nominal. Ese es un techo
duro para toda variante que le dé a la red solo el estado.

\subsection{La decisión: memoria en la entrada, no en el estado}

Hay dos formas de darle memoria a un modelo. La primera es un estado interno
propio que el modelo infiere y propaga, que es la \term{latent ODE}. Ya estaba
medida y fracasó: $R^2 = -0{,}51$, la peor de todas y la más cara. El motivo es
que ese estado latente hay que inferirlo de los datos, que es un problema mucho
más difícil que el original.

La segunda es \term{memoria en la entrada}: se le pasan a la red las últimas
$K$ muestras del comando, que son conocidas y no hay que inferir nada. En
identificación de sistemas esto es un modelo \term{NARX} o un \term{FIR} sobre
la entrada. Con la historia adentro, $\Delta f$ \emph{sí} es una función
continua de la entrada de la red, y el teorema aplica sin asteriscos.

\subsection{Una corrección al propio plan}

La primera versión fijaba $K = 100$ muestras, justificado como $5\tau$. Eso es
darle física a la red, aunque sea poca: usa el valor de $\tau$, que es justo lo
que tendría que descubrir. Se reemplazó por $K = 400$ muestras (20~ms), muy por
encima de cualquier constante de tiempo plausible, dejando que el núcleo
aprendido decaiga solo. El largo efectivo de la memoria pasa de ser un supuesto
a ser una medición.

\section{Lo que se midió}

Todas las corridas usan ventana de 100 pasos, 1500 épocas, semilla 0, y el mismo
dataset, para que sean comparables entre sí.

\subsection{Eje 1: la historia del comando}

\begin{center}
\small
\begin{tabular}{llrrrr}
\toprule
\textbf{corrida} & \textbf{qué ve $g$} & \textbf{pesos} &
\textbf{NRMSE} & \textbf{$R^2$} & \textbf{err.\ param} \\
\midrule
\texttt{e2\_wb}   & nada, sin corrección      & 0 & 15{,}23 & $-0{,}006$ & 30{,}7 \\
\texttt{e2\_B}    & $(I,E)$                   & 1\,218 & 15{,}40 & $-1{,}875$ & 34{,}4 \\
\texttt{e2\_A}    & $(I,E,P,Q)$ instantáneo   & 1\,346 & 15{,}18 & $-0{,}714$ & 32{,}9 \\
\texttt{e2\_H100} & $(I,E)$ + 100 retardos    & 7\,618 & 12{,}92 & $+0{,}250$ & 26{,}3 \\
\texttt{e2\_H400} & $(I,E)$ + 400 retardos    & 26\,818 & 11{,}48 & $+0{,}146$ & 22{,}7 \\
\rowcolor{suave}
\texttt{e3\_K400} & ídem, filtrado por un FIR & 4\,546 & \textbf{10{,}16} & $-0{,}640$ & \textbf{16{,}7} \\
\midrule
\texttt{e2\_lag2} & forma estructural, sin red & 12 & 2{,}96 & $+0{,}942$ & 5{,}1 \\
\bottomrule
\end{tabular}
\end{center}

\noindent\textbf{Lo que confirma la predicción central.} Las tres variantes sin
memoria empatan o empeoran respecto de no corregir nada. Las variantes con
memoria le ganan al white-box por primera vez, y el $R^2$ se vuelve positivo,
rompiendo el techo de $-0{,}11$ que ninguna función del estado observado podía
superar. La información faltante era la que se predijo.

\noindent\textbf{Lo que no se predijo.} Las tres métricas se mueven en
direcciones distintas al aumentar el campo receptivo. El NRMSE mejora de forma
monótona ($12{,}92 \to 11{,}48 \to 10{,}16$) y el error de parámetros también
($26{,}3 \to 22{,}7 \to 16{,}7$), pero el $R^2$ se degrada y termina negativo
($+0{,}250 \to +0{,}146 \to -0{,}640$).

La lectura habitual de un $R^2$ negativo es que la corrección absorbe física en
el lugar equivocado. Bajo esa lectura el $\beta$ identificado se tendría que
degradar, y está mejorando casi a la mitad del error. Las dos cosas no pueden
ser ciertas a la vez.

La explicación más plausible es que el $R^2$ se calcula punto a punto contra
$\Delta f$, mientras que el NRMSE mide el efecto acumulado sobre la trayectoria.
Un término que difiere punto a punto pero produce el mismo efecto dinámico
puntúa mal en el primero y bien en el segundo. La sección siguiente da evidencia
directa de por dónde va esa diferencia.

\subsection{Eje 1.3: por qué un filtro y no más entradas crudas}

La variante \texttt{H} le da a la red los $2K$ retardos crudos. Con $K=400$ son
802 entradas, y solo la primera capa densa pesa 26\,mil parámetros. Esa capa
aprende un peso independiente por cada retardo, lo cual no supone nada pero
cuesta mucho y deja mucho margen para ajustar cualquier cosa.

La variante \texttt{K} interpone un \term{FIR} (\emph{finite impulse response},
un filtro de respuesta finita al impulso) de cuatro canales antes de la red. Un
FIR calcula $\sum_k w_k\,u(t-k)$, que es una convolución con un núcleo libre.
No se le dice qué forma tiene el núcleo. Lo único que se supone es
\term{invarianza temporal}, o sea que el efecto de un retardo dado no cambia
según el instante absoluto. Eso es un supuesto de estructura, no de física.

El resultado es 4546 parámetros contra 26818, casi seis veces menos, y mejor
NRMSE. Y hay un beneficio adicional que ninguna otra variante ofrece: el núcleo
aprendido se puede graficar y comparar con la física verdadera.

\subsection{El núcleo aprendido}



\noindent Tres lecturas, en orden de importancia:

\begin{enumerate}
\item \textbf{El filtro encontró el pico en retardo cero.} La relación entre la
  cola (más allá de 10~ms) y el pico pasó de $1{,}27$ sin entrenar a $0{,}20$
  entrenado. El entrenamiento concentró energía en los retardos cortos, que es
  donde está la física. Eso es aprendizaje real, no ruido.
\item \textbf{Pero nunca suprimió la cola.} El $34{,}0\,\%$ de la energía cae en
  los primeros 5~ms, contra $25{,}2\,\%$ de un filtro sin entrenar y
  $99{,}3\,\%$ de la exponencial verdadera. El centroide temporal del núcleo
  aprendido está en $8{,}73$~ms; el de la exponencial, en $0{,}98$~ms.
\item \textbf{Eso da una explicación concreta del $R^2$ negativo.} Los pesos de
  retardo largo quedan casi libres, porque la función de costo depende muy poco
  de ellos. Su contribución a $\hat{g}$ es una componente esencialmente
  arbitraria que no se parece a $\Delta f$ y que el $R^2$ castiga punto a punto,
  mientras que el optimizador ya la ajustó para que no dañe la trayectoria.
\end{enumerate}

\begin{figure}[t]
\centering
\includegraphics[width=0.82\textwidth]{../results/figures/fir_kernel.pdf}
\caption{Energía del núcleo FIR aprendido en función del retardo, sobre
\texttt{act1}, corrida \texttt{e3\_K400} (cuatro canales, $K = 400$ muestras,
$20$~ms). La energía es la norma sobre los cuatro canales para cada retardo,
normalizada a su máximo; los canales por separado no son comparables porque la
red que sigue puede reescalar cualquiera de ellos sin cambiar la función. La
curva gris es el mismo filtro sin entrenar, y sirve de control para distinguir
estructura de ruido. La curva negra punteada es la respuesta al impulso del
actuador verdadero, $e^{-t/\tau}$ con $\tau = 1{,}0$~ms, que el entrenamiento
nunca vio.}
\end{figure}

Esto convierte una anomalía en un experimento. Si el diagnóstico es correcto,
penalizar los pesos del filtro tiene que levantar el $R^2$ sin arruinar el
NRMSE. Quedó implementado (\texttt{-{}-wd-fir}) y encolado con dos intensidades.

\section{Los errores encontrados}

Cuatro fallas de cableado aparecieron en esta tanda. Todas comparten una forma:
una lista de variantes escrita en más de un lugar del código, que se actualizó
en uno y no en los otros. Vale registrarlas porque tres de ellas no daban
síntoma.

\begin{enumerate}
\item \textbf{La variante \texttt{H} no estaba en la lista de argumentos
  válidos.} Las dos corridas murieron al instante. Es la falla barata: se ve
  enseguida.
\item \textbf{El cálculo del error open-loop no recibía el largo de la
  historia.} La red terminaba con cuatro entradas en vez de $2 + 2K$.
\item \textbf{El diagnóstico final evaluaba $g$ sobre 3000 puntos sueltos}, que
  por definición no tienen historia asociada. Se cambió para que use las propias
  ventanas de entrenamiento, que sí la tienen.
\item \textbf{El filtro FIR nunca se entrenó.} Esta es la cara. La corrida
  \texttt{e3\_K400} reportó NRMSE $20{,}90\,\%$, peor que no corregir nada,
  porque el filtro era módulo hermano de \texttt{self.g} y el optimizador arma
  su grupo con \texttt{model.g.parameters()}. Se quedó 1500 épocas en su
  inicialización aleatoria, y la red trabajó sobre cuatro proyecciones fijas y
  arbitrarias de la historia. Cuarenta y tres minutos de cómputo para un
  resultado que medía otra cosa.
\end{enumerate}

\subsection{Qué se cambió para que no se repita}

\begin{itemize}
\item \textbf{El filtro pasó a vivir dentro de \texttt{g}.} Seis lugares del
  repositorio toman \texttt{model.g.parameters()} como ``toda la corrección''.
  Hacer al filtro hijo de \texttt{g} arregla los seis a la vez, en lugar de
  agregarlo a mano en uno solo y dejar cinco mal.
\item \textbf{Un guarda en la función de entrenamiento.} Antes de construir el
  optimizador compara los parámetros entrenables del modelo contra los que
  están en sus grupos, y aborta si sobra alguno. Un parámetro fuera del
  optimizador no da síntoma: la corrida termina, reporta y guarda un número
  plausible que mide otra cosa.
\item \textbf{El checkpoint se guarda antes de evaluar.} Con el orden anterior,
  cualquier error en las últimas líneas del script costaba la corrida entera.
\item \textbf{El test de variantes chequea el gradiente por parámetro.} Antes
  sumaba sobre toda la corrección, y esa suma daba positiva igual porque la red
  sí recibía gradiente. Eso obligó a un detalle: la última capa arranca en cero
  a propósito, así que en el primer paso ninguna capa anterior recibe gradiente
  y el chequeo no distinguiría un huérfano de uno sano. Se la despierta antes de
  medir.
\item \textbf{Las decisiones de ``necesita historia'' salen de una sola
  función.} Es la causa raíz común de tres de los cuatro errores.
\end{itemize}

\section{El ruido, que hasta acá está apagado}

\subsection{Por qué no estaba}

Todo el escalado corre sobre datos limpios. Los dos datasets se generan con
\texttt{noise\_std = 0.0}, a propósito. Con ruido encima no se puede distinguir
si el residuo que la corrección tiene que aprender es el término omitido o la
medición, y el $R^2$ de la corrección dejaría de significar algo. Los
$10{,}16\,\%$ de NRMSE son sobre trayectorias exactas.

\subsection{Lo que sí está medido}

En una línea de trabajo anterior, sobre el white-box de diez parámetros. Error
máximo de parámetro, con media móvil aplicada en los cuatro niveles:

\begin{center}
\small
\begin{tabular}{lrr}
\toprule
$\sigma$ & \textbf{10 parámetros} & \textbf{subconjunto de 4} \\
\midrule
0     & 0{,}8\,\%  & 0{,}1\,\% \\
0{,}01 & 4{,}5\,\%  & 0{,}6\,\% \\
0{,}05 & 19{,}3\,\% & 1{,}2\,\% \\
0{,}10 & 41{,}3\,\% & 8{,}9\,\% \\
\bottomrule
\end{tabular}
\end{center}

La mitigación fue promediar antes de ajustar, con una media móvil de siete
muestras hasta $\sigma = 0{,}05$ y de once por encima. Lo que separa las dos
columnas es fijar los parámetros no identificables en lugar de pelearlos, y el
efecto es grande. El problema nunca fue el ruido solo, sino el ruido sobre
parámetros mal condicionados.

Hay un caso que no se mitiga y sirve de control negativo. El \term{ruido de
proceso}, el que entra dentro de la ecuación diferencial en vez de sumarse a la
medición, da $R^2$ del oráculo de $0{,}12$ con el estado y $0{,}11$ agregando el
comando. No es función de nada, así que ninguna corrección lo puede aprender, ni
con historia ni con capacidad.

\subsection{La hipótesis que se va a probar}

La variante \texttt{H} le da a la red 802 entradas y 26\,mil pesos en la primera
capa, que es exactamente la configuración que ajusta ruido. El FIR promedia $K$
muestras, y eso es un filtro pasabajos por construcción. Si el FIR se degrada
menos al subir el ruido, el argumento a favor de la convolución deja de ser solo
el costo.

Se prepararon \texttt{act1\_n01} ($\sigma = 0{,}01$, el $5{,}7\,\%$ del desvío de
$E$) y \texttt{act1\_n05} ($\sigma = 0{,}05$, el $27{,}5\,\%$). El ruido se suma
al dataset ya generado en vez de regenerar las trayectorias, por dos motivos: la
trayectoria queda idéntica a la limpia, así que las corridas son comparables
muestra a muestra, y el $\Delta f$ de referencia sigue calculado sobre la
trayectoria exacta, que es lo único que hace válido el $R^2$.

\section{Qué quedó corriendo}

Doce corridas encoladas en secuencia, nunca en paralelo. La bitácora ya
registra que cuatro carriles sobre ocho núcleos convirtieron corridas de 25
minutos en corridas de 25 horas. El lote entero son unas ocho horas y media de
reloj, contando desde el mediodía del 11 de septiembre.

\begin{center}
\small
\begin{tabular}{llp{7.4cm}}
\toprule
\textbf{corrida} & \textbf{planta} & \textbf{qué contesta} \\
\midrule
\texttt{e3\_K800} & \texttt{act1} & Si el campo receptivo más largo sigue mejorando. \\
\texttt{e3\_K400\_h64} & \texttt{act1} & Si falta capacidad en la red posterior al filtro. \\
\texttt{e3\_K400\_f8} & \texttt{act1} & Si el filtro de cuatro canales es el cuello de botella. Con ocho canales, si la hipótesis del rango es correcta, el $R^2$ tiene que recuperarse. \\
\texttt{e4\_Sg} & \texttt{refrac1} & Cuánto se tapan entre sí los parámetros y la red, sobre la planta donde la forma exacta está completa. Si la red no aprende nada, su magnitud queda cerca de cero. \\
\texttt{e4\_S2ver} & \texttt{refrac1} & Verificación del dataset. Tiene que reproducir NRMSE $1{,}70$ y $R^2$ $0{,}964$. \\
\texttt{e4\_wb\_n05}, \texttt{e4\_K400\_n05}, \texttt{e4\_H400\_n05} & \texttt{act1\_n05} & El eje 6 con $\sigma = 0{,}05$, con el white-box como referencia. \\
\texttt{e4\_K400\_n01}, \texttt{e4\_H400\_n01} & \texttt{act1\_n01} & Ídem con $\sigma = 0{,}01$. \\
\texttt{e5\_K400\_wd1e-4}, \texttt{e5\_K400\_wd1e-3} & \texttt{act1} & Si penalizar los pesos del filtro levanta el $R^2$ sin arruinar el NRMSE. \\
\bottomrule
\end{tabular}
\end{center}

\section{Dónde está el proyecto}

La escala honesta: la forma estructural sigue ganando por lejos, $2{,}96$ contra
$10{,}16$ en NRMSE y $0{,}942$ contra $-0{,}640$ en $R^2$. El camino
generalizable funciona y rompió el techo teórico, pero todavía está a un tercio
del camino en reproducción. Escribir la ecuación a mano sigue siendo mucho mejor
que aprenderla.

Lo que cambió respecto del 10 de septiembre es que ahora hay un camino que
mejora al agrandarlo, y una explicación medible de por qué todavía no llega.
Antes del eje 1 no había ninguno de los dos.

\subsection{Lo que queda decidido}

La prioridad pasó a ser el NRMSE. El $R^2$ de la corrección se reporta siempre y
no bloquea nada. El criterio de aceptación es NRMSE por debajo de $10{,}16\,\%$,
con el white-box ($15{,}23$) como piso y la forma estructural ($2{,}96$) como
referencia de cuánto falta.

El control que mantiene honesta esa degradación es el error de parámetros. Si la
corrección mejorara el NRMSE absorbiendo física en el lugar equivocado, el
$\beta$ identificado se degradaría. Está mejorando. Mientras las dos cosas se
muevan juntas, la mejora es real y no un canje.

\subsection{Lo que queda pendiente de una decisión}

\begin{itemize}
\item \textbf{Los datos reales.} Para probar cualquier variante sobre una
  grabación hay que estimar el estado inicial de cada ventana, porque lo que se
  observa no es $(I,E)$ sino una proyección escalar $s = c\,(E - I)$. La
  recomendación es optimizar solo el estado inicial con los parámetros
  congelados, pero es una decisión de método y no se tomó.
\item \textbf{La revisión del repositorio limpio} antes de hacerlo público. Ya
  aparecieron tres afirmaciones falsas en su README, así que conviene asumir que
  hay más.
\end{itemize}

\end{document}
